% Alternative main file for Peking Mathematical Journal submission
% Uses standard article class (if svjour3 is not available)
% Compile with: pdflatex pmj_submission_article.tex

\documentclass[11pt,a4paper]{article}

%=============================================================================
% PAGE LAYOUT (Springer-compatible)
%=============================================================================
\usepackage[
  top=2.5cm,
  bottom=2.5cm,
  left=2.5cm,
  right=2.5cm
]{geometry}

%=============================================================================
% PACKAGES
%=============================================================================
\usepackage{amsmath,amssymb,amsthm}
\usepackage{mathtools}
\usepackage[T1]{fontenc}
\usepackage{hyperref}
\usepackage{cleveref}
\usepackage{tikz}
\usetikzlibrary{positioning}
\usepackage{booktabs}
\usepackage{enumitem}
\usepackage[hypcap=false]{caption}
\usepackage{graphicx}
\usepackage{authblk}  % For author affiliations

% For review mode, uncomment:
% \usepackage{lineno}
% \linenumbers
% \usepackage{setspace}
% \doublespacing

%=============================================================================
% HYPERREF SETUP
%=============================================================================
\hypersetup{
  colorlinks=true,
  linkcolor=blue,
  citecolor=blue,
  urlcolor=blue,
  pdftitle={The Angular Momentum Penrose Inequality},
  pdfauthor={Da Xu},
  pdfkeywords={Penrose inequality, angular momentum, Kerr spacetime, MOTS}
}

%=============================================================================
% THEOREM ENVIRONMENTS
%=============================================================================
\theoremstyle{plain}
\newtheorem{theorem}{Theorem}[section]
\newtheorem{conjecture}[theorem]{Conjecture}
\newtheorem{proposition}[theorem]{Proposition}
\newtheorem{lemma}[theorem]{Lemma}
\newtheorem{corollary}[theorem]{Corollary}

\theoremstyle{definition}
\newtheorem{definition}[theorem]{Definition}
\newtheorem{example}[theorem]{Example}

\theoremstyle{remark}
\newtheorem{remark}[theorem]{Remark}

%=============================================================================
% CUSTOM COMMANDS
%=============================================================================
\newcommand{\ADM}{\mathrm{ADM}}
\newcommand{\R}{\mathbb{R}}
\newcommand{\MOTS}{\mathrm{MOTS}}
\newcommand{\DEC}{\mathrm{DEC}}
\newcommand{\tr}{\mathrm{tr}}
\newcommand{\Ric}{\mathrm{Ric}}
\newcommand{\Rm}{\mathrm{Rm}}
\newcommand{\Div}{\mathrm{div}}
\DeclareMathOperator{\ind}{ind}
\DeclareMathOperator{\coker}{coker}
\newcommand{\tM}{\tilde{M}}
\newcommand{\tg}{\tilde{g}}
\newcommand{\bg}{\bar{g}}
\newcommand{\bM}{\bar{M}}
\newcommand{\momdens}{\boldsymbol{j}}
\newcommand{\Hoelder}{\beta}
\newcommand{\Komarform}{\alpha_J}

% Allow page breaks in equations
\allowdisplaybreaks

% Mild help for line breaking
\emergencystretch=1em
\tolerance=2000

\setcounter{tocdepth}{2}

%=============================================================================
\begin{document}

%=============================================================================
% TITLE PAGE
%=============================================================================
\title{\textbf{The Angular Momentum Penrose Inequality}}

\author{Da Xu}
\affil{China Mobile Research Institute, Beijing 100053, China\\
\texttt{xuda@chinamobile.com}}

\date{}

\maketitle

%=============================================================================
% ABSTRACT
%=============================================================================
\begin{abstract}
\noindent
We establish the Penrose inequality with angular momentum for asymptotically flat, axisymmetric vacuum initial data containing a strictly stable marginally outer trapped surface. Specifically, if $(M^3, g, K)$ satisfies the dominant energy condition and $\Sigma$ denotes the outermost MOTS with area $A$ and Komar angular momentum $J$, then
\[
M_{\ADM} \geq \sqrt{\frac{A}{16\pi} + \frac{4\pi J^2}{A}},
\]
with equality characterizing slices of Kerr. The argument combines the Jang equation approach of Bray--Khuri and Han--Khuri with the $p$-harmonic level set method of Agostiniani--Mazzieri--Oronzio. A central role is played by a modified Hawking mass that incorporates angular momentum and remains monotone along the flow. The conservation of $J$ under this flow relies on the vacuum hypothesis through a cohomological argument. We also characterize the equality case using the Mars--Simon tensor.

\medskip
\noindent
\textbf{Keywords:} Penrose inequality, angular momentum, Kerr spacetime, marginally outer trapped surface, Jang equation, conformal method

\medskip
\noindent
\textbf{Mathematics Subject Classification (2020):} 83C57, 53C21, 35J60, 83C40
\end{abstract}

\tableofcontents
\newpage

%=============================================================================
% MAIN CONTENT
%=============================================================================

% Section 1: Introduction
\input{sec01-introduction}

% Section 2: Verification for Kerr Spacetime
\section{Verification for Kerr Spacetime}\label{sec:kerr}
%=============================================================================

Before proceeding to the proof, we verify that Kerr spacetime saturates the inequality. This is a necessary consistency check.

Working in geometric units ($G=c=1$), the Kerr spacetime with mass $M$ and spin parameter $a = J/M$ (where $|a| \leq M$) has outer horizon radius $r_+ = M + \sqrt{M^2 - a^2}$ and horizon area
\[
A = 4\pi(r_+^2 + a^2) = 8\pi M(M + \sqrt{M^2 - a^2}).
\]

\begin{theorem}[Kerr Saturation]\label{thm:kerr}
For all sub-extremal Kerr black holes ($|a| \leq M$),
\[
M = \sqrt{\frac{A}{16\pi} + \frac{4\pi J^2}{A}}.
\]
\end{theorem}

\begin{proof}
Set $s = \sqrt{M^2 - a^2}$, so $r_+ = M + s$. Then
\[
\frac{A}{16\pi} = \frac{M(M + s)}{2}, \qquad \frac{4\pi J^2}{A} = \frac{Ma^2}{2(M + s)}.
\]
Adding these gives
\[
\frac{A}{16\pi} + \frac{4\pi J^2}{A} = \frac{M[(M + s)^2 + a^2]}{2(M + s)}.
\]
Since $s^2 = M^2 - a^2$, we have $(M + s)^2 + a^2 = 2M(M + s)$, and so
\[
\frac{A}{16\pi} + \frac{4\pi J^2}{A} = M^2. \qedhere
\]
\end{proof}

As special cases: for Schwarzschild ($a = 0$), we have $A = 16\pi M^2$ and $J = 0$, giving the standard bound; for extremal Kerr ($a = M$), we have $A = 8\pi M^2$ and $J = M^2$, and the two terms contribute equally to give $M$.

\begin{remark}[Kerr Data Regularity]\label{rem:kerr-regularity}
On a Boyer--Lindquist slice with $r > r_+$, Kerr initial data satisfies $g_{ij} - \delta_{ij} = O(M/r)$ and $K_{ij} = O(Ma/r^2)$, with decay rate $\tau = 1 > 1/2$. The data satisfies hypotheses (H1)--(H4) for strictly sub-extremal parameters; see \cite{ChruscielCostaHeusler2012}.
\end{remark}

%=============================================================================


% Section 3: Proof Strategy Overview
\input{sec03-proof-outline}

% Section 4: Stage 1 - Axisymmetric Jang Equation
\input{sec04-jang}

% Section 5: Stage 2 - AM-Lichnerowicz Equation
\input{sec05-lichnerowicz}

% Section 6: Stage 3 - AMO Flow with Angular Momentum
\input{sec06-amo}

% Section 7: Stage 4 - Sub-Extremality
\input{sec07-subextremality}

% Section 8: Synthesis - Complete Proof
\input{sec08-synthesis}

% Section 9: Rigidity
\input{sec09-rigidity}

% Section 10: Extensions and Open Problems
\input{sec10-extensions}

% Section 11: Conclusion
\input{sec13-conclusion}

%=============================================================================
\appendix
%=============================================================================

% Appendix A: Numerical Illustrations
\input{sec11-numerical}

% Appendix B: Technical Foundations
\input{sec12-technical}

% Appendix C: Key AMO Estimates
\input{app01-amo-estimates}

% Appendix D: Schauder Estimates
\input{app02-schauder}

% Appendix E: Supersolution Analysis
\input{app03-supersolution}

% Appendix F: Sub-Extremality Factor Improvement
\input{app04-subext-improvement}

% Appendix G: Mars-Simon Tensor
\input{app05-mars-simon}

% Appendix H: Function Space Compatibility Verification
\input{app06-function-spaces}

%=============================================================================
% ACKNOWLEDGMENTS
%=============================================================================
\section*{Acknowledgments}
The author would like to thank China Mobile Research Institute for supporting fundamental research.

%=============================================================================
% DECLARATIONS (Required by Springer journals)
%=============================================================================
\section*{Declarations}

\paragraph{Conflict of Interest.} The author has no conflicts of interest to declare that are relevant to the content of this article.

\paragraph{Data Availability.} No datasets were generated or analyzed during this study.

%=============================================================================
% BIBLIOGRAPHY
%=============================================================================
\bibliographystyle{plain}
\bibliography{references}

\end{document}
